\chapter{Complexity — NARM Logspace}
%%% generated by the blueprint pipeline from TCSlib.Complexity.SpaceComplexity.Machines.NARMLogspace
\section{Overview}

Declarations of the Lean module \texttt{TCSlib.Complexity.SpaceComplexity.Machines.NARMLogspace}.

\section{Declarations}

\begin{theorem}[Logspace Simulation Preserves Output Probabilities]
\lean{Complexity.LogProg.narm_logspace_ptm}
\label{Complexity.LogProg.narm_logspace_ptm}
\leanok
\uses{Complexity.LOGSPACE, Complexity.LogProg.AConf, Complexity.LogProg.NARM, Complexity.LogProg.NGuard, Complexity.LogProg.NHalts, Complexity.LogProg.NHalts.mono, Complexity.LogProg.NPreS, Complexity.LogProg.ansProb, Complexity.LogProg.bankK, Complexity.LogProg.bankStart, Complexity.LogProg.bankTM, Complexity.LogProg.length_bits_le_log, Complexity.LogProg.log_poly_bound, Complexity.LogProg.narm_space_randProb, Complexity.LogProg.ncompileFinNDTM, Complexity.LogProg.npre_of_npreS, Complexity.logSpace, Turing.FinNDTM, Turing.FinNDTM.UsesSpace, Turing.FinNDTM.outputProb, Turing.FinNDTM.outputWith, Turing.FinNDTM.usesSpace_of_haltsWithin, Turing.FinTM, Turing.FinTM.DecidesInSpace, Turing.MultiTapeTM.indicator, Turing.NDTM.HaltsWithin, Turing.NDTM.initCfg, Turing.NDTM.spaceUsedWith}
\difficulty{6}
Let a nondeterministic register program with $m$ registers and $d$ oracle languages have
a finite control-state set with decidable equality, and suppose every oracle language
belongs to $\text{LOGSPACE}$. For natural numbers $C_0,c_0$, assume that for every
input $x$ some length $N$ makes the program halt on every choice string of length $N$,
with each pre-final configuration satisfying its validity conditions and having every
register value at most $C_0(|x|+1)^{c_0}$. Then there are a natural number $c$ and a
finite-state nondeterministic machine that uses at most $c$ times logarithmic space on
every input, and such that for every $x,N$ for which the stated halting and
configuration conditions hold, the machine halts within some time $T$ and, for each bit
$b$, the probability that it outputs the one-bit string $[b]$ after $T$ steps equals the
probability that the register program answers $b$ after $N$ choices.
\end{theorem}

\begin{theorem}[Register Programs in Randomized Logarithmic Space]
\lean{Complexity.LogProg.narm_mem_RL}
\label{Complexity.LogProg.narm_mem_RL}
\leanok
\uses{Complexity.LOGSPACE, Complexity.LogProg.NARM, Complexity.LogProg.NHalts, Complexity.LogProg.NPreS, Complexity.LogProg.ansProb, Complexity.LogProg.narm_logspace_ptm, Randomized.RL, Randomized.exists_ofFn_eq, Randomized.forall_of_randProb_eq_one, Randomized.mem_RSPACE_of_pos, Turing.MultiTapeTM.indicator}
\difficulty{4}
Let $L$ be a language of binary strings, and let a nondeterministic register program
with $m$ registers and $d$ oracle languages have a finite control-state set with
decidable equality. Suppose every oracle language belongs to $\text{LOGSPACE}$, and let
$C_0,c_0$ be natural numbers and $\delta$ a positive rational. Assume that for each
input $x$ there is a length $N$ such that every choice string of length $N$ makes the
program halt, every pre-final configuration satisfies its validity conditions and has
each register value at most $C_0(|x|+1)^{c_0}$, and the following probability conditions
hold: if $x\in L$, the probability of answering yes is at least $\delta$; if $x\notin
L$, the probability of answering no is $1$. Then $L$ belongs to $\text{RL}$, the class
decidable by a randomized machine using logarithmic space with one-sided error.
\end{theorem}
